\documentclass[review,11pt,authoryear]{elsarticle}
\usepackage[margin=2.5cm]{geometry}

\usepackage{lineno}
\usepackage{graphicx}
\usepackage{booktabs}
\usepackage{amsmath}
\usepackage{amssymb}
\usepackage{hyperref}
\usepackage{multirow}
\usepackage{xcolor}
\usepackage{listings}
\usepackage{placeins}

\modulolinenumbers[1]

% Code listing style
\lstdefinestyle{python}{
  language=Python,
  basicstyle=\ttfamily\small,
  keywordstyle=\color{blue},
  stringstyle=\color{red!70!black},
  commentstyle=\color{green!50!black},
  numbers=none,
  frame=single,
  breaklines=true,
  showstringspaces=false,
  tabsize=4,
  xleftmargin=2em,
  framexleftmargin=1.5em,
  aboveskip=1em,
  belowskip=1em
}

\lstdefinestyle{bash}{
  language=bash,
  basicstyle=\ttfamily\small,
  keywordstyle=\color{blue},
  stringstyle=\color{red!70!black},
  commentstyle=\color{green!50!black},
  numbers=none,
  frame=single,
  breaklines=true,
  showstringspaces=false,
  tabsize=4,
  xleftmargin=2em,
  framexleftmargin=1.5em,
  aboveskip=1em,
  belowskip=1em
}

\lstdefinestyle{yaml}{
  basicstyle=\ttfamily\small,
  commentstyle=\color{green!50!black},
  numbers=none,
  frame=single,
  breaklines=true,
  showstringspaces=false,
  tabsize=2,
  xleftmargin=2em,
  framexleftmargin=1.5em,
  aboveskip=1em,
  belowskip=1em,
  morecomment=[l]{\#}
}

\journal{MethodsX}

\begin{document}

\begin{frontmatter}

\title{Remote SAMsing: A Python Pipeline for SAM2-Based Segmentation of Large Remote Sensing Images}

\author[ee]{Osmar Luiz Ferreira de Carvalho}
\ead{osmarcarvalho@ieee.org}

\author[geo]{Osmar Ab\'{i}lio de Carvalho J\'{u}nior\corref{cor1}}
\ead{osmarjr@unb.br}

\author[geo]{Anesmar Olino de Albuquerque}
\ead{anesmar@ieee.org}

\author[ee]{Daniel Guerreiro e Silva}
\ead{danielgs@unb.br}

\cortext[cor1]{Corresponding author}
\address[ee]{Department of Electrical Engineering, University of Bras\'{i}lia, Bras\'{i}lia, Brazil}
\address[geo]{Department of Geography, University of Bras\'{i}lia, Bras\'{i}lia, Brazil}

\begin{abstract}
\begin{itemize}
\item Remote SAMsing is an open-source Python pipeline that segments arbitrarily large remote sensing images using SAM2 (Segment Anything Model 2) without modifying the model's architecture or weights.
\item The pipeline addresses three limitations of applying SAM2 directly to remote sensing imagery: incomplete coverage from single-pass inference, boundary fragmentation from tiling, and the absence of scale-aware configuration guidance.
\item Remote SAMsing is available as a pip-installable package with both a command-line interface and a Python API.
\end{itemize}

The Segment Anything Model 2 (SAM2) produces high-quality segmentation masks on natural images, but a single inference pass on remote sensing imagery leaves 30--70\% of the scene unsegmented, and the model's fixed $1{,}024 \times 1{,}024$ input resolution prevents direct application to images spanning tens of thousands of pixels. Remote SAMsing solves both problems through a multi-pass algorithm that iteratively focuses SAM2 on residual unsegmented areas using black mask focusing and adaptive threshold decay, combined with a parameter-free boundary merge based on Union-Find that produces spatially consistent label mosaics with constant GPU memory.

The pipeline accepts any raster image readable by GDAL/rasterio and produces georeferenced outputs: a labeled GeoTIFF raster (each pixel assigned a segment ID), vectorized polygons (Shapefile and/or GeoPackage), and processing statistics in JSON format. All outputs preserve the coordinate reference system and spatial transform of the input. The pipeline runs on any NVIDIA GPU with 8~GB or more of VRAM and processes images of arbitrary size through tile-by-tile streaming.

Remote SAMsing exposes a single main configuration parameter---tile size---that functions as an implicit scale parameter: smaller tiles magnify objects in SAM2's internal resolution, improving detection of small features such as cars and narrow walkways. The companion paper \citep{carvalho2026jag} provides a full scientific evaluation on seven scenes spanning 5~cm to 4.78~m ground sampling distance, demonstrating 91--98\% coverage, Det@0.5 of 62--96\%, and Boundary IoU 3--8$\times$ higher than SLIC and Felzenszwalb baselines. A scalability test on a $36{,}000 \times 54{,}000$ pixel mosaic (1.94 billion pixels) produced 124{,}180 segments at 97\% coverage.

Source code: \url{https://github.com/osmarluiz/sam-mosaic}
\end{abstract}

\begin{keyword}
SAM2 \sep remote sensing \sep image segmentation \sep superpixels \sep tiling \sep OBIA \sep foundation models \sep Python
\end{keyword}

\end{frontmatter}

\linenumbers


% ============================================================================
\section{Method Details}
\label{sec:method}
% ============================================================================

% ----------------------------------------------------------------------------
\subsection{Installation and Requirements}
\label{sec:install}
% ----------------------------------------------------------------------------

Remote SAMsing requires Python 3.12 or later, PyTorch 2.x with CUDA support, and an NVIDIA GPU with at least 8~GB of VRAM. The tested configuration uses PyTorch 2.9, CUDA 12.8, and the SAM2.1 Hiera Large checkpoint (${\sim}857$~MB).

\subsubsection{Installation from source}

\begin{lstlisting}[style=bash]
# Clone the repository
git clone https://github.com/osmarluiz/sam-mosaic.git
cd sam-mosaic

# Install the package
pip install -e .

# Install SAM2 from PyPI
pip install sam2
\end{lstlisting}

\noindent
The \texttt{sam2} package must be installed from PyPI (JinsuaFeito-dev fork, version 1.1.0+), which has been tested for stable GPU memory during large-scale processing. The official Facebook repository version may cause memory leaks during extended runs.

\subsubsection{Checkpoint download}

\begin{lstlisting}[style=bash]
cd checkpoints
wget https://dl.fbaipublicfiles.com/segment_anything_2/092824/sam2.1_hiera_large.pt
\end{lstlisting}

\noindent
The checkpoint is the standard SAM2.1 Hiera Large model distributed by Meta AI. Remote SAMsing does not modify the model weights.

\subsubsection{Conda environment (alternative)}

For exact reproducibility, a conda environment file is provided:

\begin{lstlisting}[style=bash]
conda env create -f environment.yml
conda activate ts_annotator
\end{lstlisting}

\subsubsection{Hardware requirements}

Table~\ref{tab:hardware} lists hardware requirements. The streaming mode (Section~\ref{sec:large}) allows processing images larger than available RAM by writing intermediate results to disk.

\begin{table}[htbp]
\centering
\caption{Hardware requirements for Remote SAMsing.}
\label{tab:hardware}
\begin{tabular}{lll}
\toprule
Component & Minimum & Tested \\
\midrule
GPU VRAM & 8 GB & 24 GB (RTX 4090) \\
System RAM & 16 GB & 64 GB \\
Python & 3.12+ & 3.12 \\
PyTorch & 2.0+ & 2.9 \\
CUDA & 12.0+ & 12.8 \\
Disk & 1 GB + output & SSD recommended \\
\bottomrule
\end{tabular}
\end{table}


% ----------------------------------------------------------------------------
\subsection{Pipeline Overview}
\label{sec:overview}
% ----------------------------------------------------------------------------

Remote SAMsing processes large images in four stages:

\begin{enumerate}
    \item \textbf{Tiling.} The input image is divided into a regular grid of $T \times T$ pixel tiles. Each tile is extracted with $p$ pixels of contextual padding on each side, producing a $(T + 2p) \times (T + 2p)$ window for inference. Default: $T = 1{,}000$, $p = 50$.

    \item \textbf{Multi-pass segmentation.} Each tile undergoes iterative SAM2 inference. The first pass runs SAM2's Automatic Mask Generator (AMG) with a uniform $64 \times 64$ point grid at high quality thresholds ($\tau = 0.93$). Subsequent passes apply three mechanisms: (a) \textit{black mask focusing}---already-segmented pixels are set to black, preventing SAM2 from re-segmenting them; (b) \textit{dense grid sampling}---the same uniform grid is regenerated and points falling on segmented pixels are discarded; (c) \textit{adaptive threshold decay}---when coverage gain stagnates (below 0.1 percentage points), both the IoU and stability thresholds decrease by a step of 0.01. The loop stops when coverage reaches the target (99\%), thresholds are exhausted ($\tau_{\text{end}} = 0.60$), or the maximum number of passes is reached.

    \item \textbf{Boundary merge.} After all tiles are segmented, a best-match merge algorithm reconciles labels at tile boundaries. For each boundary, the algorithm records contact counts between label pairs on adjacent pixels. Each segment selects its best merge partner (the neighbor with the highest contact count). A Union-Find structure with path compression resolves transitive merges, and a lookup table relabels the entire mosaic in a single pass. The operation is parameter-free and runs in effectively linear time.

    \item \textbf{Post-processing and output.} Small enclosed fragments are absorbed into their surrounding segment. Isolated components below a minimum area threshold ($a_{\min} = 100$ pixels) are removed. The labeled raster is saved as a georeferenced GeoTIFF, and polygons are vectorized with optional Douglas-Peucker simplification.
\end{enumerate}

\noindent
The companion paper \citep{carvalho2026jag} describes the algorithm in full detail, including formal pseudocode and ablation experiments that quantify the contribution of each component.


% ----------------------------------------------------------------------------
\subsection{Python API}
\label{sec:api}
% ----------------------------------------------------------------------------

Remote SAMsing provides two Python functions as entry points: \texttt{segment\_image()} for standard usage with configuration files or defaults, and \texttt{segment\_with\_params()} for direct parameter control.

\subsubsection{Minimal usage (3 lines)}

\begin{lstlisting}[style=python]
from sam_mosaic import segment_image

result = segment_image(
    "input.tif",
    "output/",
    checkpoint="checkpoints/sam2.1_hiera_large.pt"
)

print(f"Segments: {result.n_segments}")
print(f"Coverage: {result.coverage:.1f}%")
print(f"Time: {result.processing_time:.1f}s")
\end{lstlisting}

\noindent
This runs the full pipeline with default parameters (Table~\ref{tab:params}). The \texttt{SegmentationResult} object contains output file paths (\texttt{labels\_path}, \texttt{shapefile\_path}) and statistics (\texttt{n\_segments}, \texttt{coverage}, \texttt{processing\_time}).

\subsubsection{Urban scene with small objects}

For scenes with small objects (cars, narrow streets) at resolutions finer than 50~cm, reducing tile size improves detection:

\begin{lstlisting}[style=python]
from sam_mosaic import segment_with_params

result = segment_with_params(
    "urban_24cm.tif",
    "output/urban/",
    checkpoint="checkpoints/sam2.1_hiera_large.pt",
    tile_size=250,           # 4x magnification
    point_strategy="dense_grid",
    erosion_iterations=0,
)
\end{lstlisting}

\subsubsection{Agricultural scene with large fields}

For coarser-resolution imagery with large, spectrally homogeneous objects, the default tile size of 1000 produces the best boundary quality:

\begin{lstlisting}[style=python]
result = segment_with_params(
    "agriculture_5m.tif",
    "output/agriculture/",
    checkpoint="checkpoints/sam2.1_hiera_large.pt",
    tile_size=1000,          # default
    point_strategy="dense_grid",
)
\end{lstlisting}

\subsubsection{Ablation study}

The \texttt{segment\_with\_params()} function exposes all internal parameters for reproducible experiments:

\begin{lstlisting}[style=python]
# Single-pass only (no multi-pass)
result = segment_with_params(
    "input.tif", "output/single_pass/",
    checkpoint="checkpoints/sam2.1_hiera_large.pt",
    max_passes=1,
    iou_start=0.86,
    use_adaptive_threshold=False
)

# Without black mask
result = segment_with_params(
    "input.tif", "output/no_blackmask/",
    checkpoint="checkpoints/sam2.1_hiera_large.pt",
    use_black_mask=False
)

# Without contextual padding (to study merge artifacts)
result = segment_with_params(
    "input.tif", "output/no_padding/",
    checkpoint="checkpoints/sam2.1_hiera_large.pt",
    padding=0
)
\end{lstlisting}

\subsubsection{Configuration-file usage}

The \texttt{segment\_image()} function also accepts a path to a YAML configuration file:

\begin{lstlisting}[style=python]
from sam_mosaic import segment_image

result = segment_image(
    "input.tif",
    "output/",
    config="my_config.yaml"
)
\end{lstlisting}


% ----------------------------------------------------------------------------
\subsection{Configuration Parameters}
\label{sec:params}
% ----------------------------------------------------------------------------

Table~\ref{tab:params} lists all configurable parameters with their default values, valid ranges, and usage guidance.

\begin{table}[htbp]
\centering
\caption{Complete parameter reference for Remote SAMsing. Parameters are organized by functional group. Default values reflect the recommended configuration for general-purpose remote sensing segmentation.}
\label{tab:params}
\scriptsize
\begin{tabular}{p{3.8cm}lp{6cm}}
\toprule
Parameter & Default & Description and guidance \\
\midrule
\multicolumn{3}{c}{\textit{Tiling}} \\
\midrule
\texttt{tile\_size} & 1000 & Tile side length in pixels. Functions as an implicit scale parameter: smaller values magnify objects in SAM2's $1{,}024^2$ internal resolution. Use 250--500 for urban scenes with small objects at $<$50~cm GSD; use 1000 for agricultural scenes with large fields. \\
\texttt{padding} & 50 & Contextual padding in pixels per side. Ensures segments extend to tile edges for merge. Increase to 100 if merge artifacts appear at tile boundaries. \\
\midrule
\multicolumn{3}{c}{\textit{Thresholds}} \\
\midrule
\texttt{iou\_start} & 0.93 & Initial predicted IoU threshold. Higher values accept only high-confidence masks in early passes. \\
\texttt{iou\_end} & 0.60 & Final IoU threshold. The loop stops when thresholds reach this value. Lower values allow more permissive masks in late passes. \\
\texttt{stability\_start} & 0.93 & Initial stability score threshold. Mirrors \texttt{iou\_start} for the stability filter. \\
\texttt{stability\_end} & 0.60 & Final stability threshold. Mirrors \texttt{iou\_end}. \\
\texttt{threshold\_step} & 0.01 & Threshold decrease per stagnation event. \\
\midrule
\multicolumn{3}{c}{\textit{Segmentation}} \\
\midrule
\texttt{points\_per\_side} & 64 & Grid density: $64 \times 64 = 4{,}096$ prompt points per tile. Higher values increase point density but also processing time. \\
\texttt{target\_coverage} & 99.0 & Coverage percentage at which the multi-pass loop stops. \\
\texttt{max\_passes} & None & Maximum passes per tile. None means unlimited (loop stops by coverage or threshold exhaustion). \\
\texttt{use\_black\_mask} & True & Mask already-segmented pixels with black between passes. Disabling reduces coverage by 4--11 pp. \\
\texttt{use\_adaptive\_threshold} & True & Decrease thresholds when progress stagnates. Disabling reduces coverage by 3--20 pp. \\
\texttt{point\_strategy} & kmeans & Point selection in subsequent passes: \texttt{"kmeans"} clusters points in residual areas; \texttt{"dense\_grid"} regenerates a uniform grid filtered by the segmentation mask. Dense Grid is recommended for most use cases. \\
\texttt{erosion\_iterations} & 0 & Morphological erosion iterations on the residual mask before placing points. Higher values keep points away from segment edges. \\
\texttt{crop\_n\_layers} & 0 & SAM2 multi-scale crop layers. 0 = single scale. 1 = original + 4 sub-crops at $2\times$ zoom. Increases processing time 3--5$\times$; tile size reduction is more effective for small objects. \\
\texttt{box\_nms\_thresh} & 0.7 & NMS threshold for overlapping mask bounding boxes. 1.0 disables NMS entirely. \\
\texttt{sort\_by\_area} & False & Sort masks by area ascending before applying. Gives small objects (cars) priority over large ones (rooftops). \\
\midrule
\multicolumn{3}{c}{\textit{Merge}} \\
\midrule
\texttt{merge\_strategy} & best\_match & Boundary merge strategy: \texttt{"best\_match"} (parameter-free, recommended), \texttt{"mutual\_best"} (stricter), or \texttt{"min\_contact"} (threshold-based). \\
\texttt{min\_contact\_pixels} & 20 & Minimum contact pixels for \texttt{min\_contact} strategy only. \\
\texttt{min\_mask\_area} & 100 & Remove segments smaller than this area (pixels) after merge. \\
\texttt{merge\_enclosed\_max\_area} & 500 & Maximum area for absorbing enclosed fragments into their surrounding segment. \\
\midrule
\multicolumn{3}{c}{\textit{Output}} \\
\midrule
\texttt{save\_labels} & True & Save labeled raster as GeoTIFF (\texttt{labels.tif}). \\
\texttt{save\_shapefile} & True & Save vectorized polygons as Shapefile (\texttt{segments.shp}). \\
\texttt{save\_geopackage} & False & Save vectorized polygons as GeoPackage (\texttt{segments.gpkg}). \\
\texttt{simplify\_tolerance} & 1.0 & Douglas-Peucker simplification tolerance in map units. 0 = no simplification. Higher values produce smaller files with less boundary detail. \\
\texttt{streaming\_mode} & auto & Mosaic storage mode: \texttt{"auto"} uses disk for large images ($>$30\% of RAM), \texttt{"ram"} keeps the mosaic in memory, \texttt{"disk"} always streams to disk. \\
\bottomrule
\end{tabular}
\end{table}

\FloatBarrier

\subsubsection{When to change tile size}

Tile size is the single most important configuration parameter. It controls the effective spatial resolution at which SAM2 processes the image, because SAM2 internally resizes all inputs to $1{,}024 \times 1{,}024$ pixels regardless of the input size. A tile of $T = 1{,}000$ pixels undergoes almost no resize; a tile of $T = 250$ is magnified $4\times$, making each pixel four times larger in SAM2's internal representation.

The recommended approach:
\begin{itemize}
    \item Start with $T = 1{,}000$ (default).
    \item If small objects (cars, narrow walkways) are not detected, reduce to $T = 500$ or $T = 250$.
    \item If large, homogeneous objects (agricultural fields, water bodies) fragment into chains of segments, increase $T$ back to 1000.
    \item As a rule of thumb, each tile should cover at least $50 \times 50$~m of ground area to contain enough scene diversity for SAM2 to generate meaningful segments.
\end{itemize}

\subsubsection{When to change thresholds}

The default threshold range ($0.93 \to 0.60$) was selected to balance quality and coverage across a range of scene types \citep{carvalho2026jag}. Adjustments are rarely needed, but two scenarios may benefit:
\begin{itemize}
    \item \textbf{Higher quality, lower coverage:} Set \texttt{iou\_end} and \texttt{stability\_end} to 0.70 or higher. This terminates the multi-pass loop earlier, producing fewer but higher-confidence segments.
    \item \textbf{Maximum coverage at the expense of quality:} Set \texttt{iou\_end} and \texttt{stability\_end} to 0.50. Late passes will accept noisier masks to fill remaining gaps.
\end{itemize}


% ----------------------------------------------------------------------------
\subsection{Output Formats}
\label{sec:output}
% ----------------------------------------------------------------------------

Remote SAMsing produces four output files in the specified output directory:

\begin{table}[htbp]
\centering
\caption{Output files generated by Remote SAMsing.}
\label{tab:outputs}
\begin{tabular}{llp{7cm}}
\toprule
File & Format & Description \\
\midrule
\texttt{labels.tif} & GeoTIFF & Labeled raster where each pixel contains an integer segment ID (1, 2, 3, \ldots). Background (unsegmented) pixels have value 0. Preserves the CRS, spatial transform, and resolution of the input image. \\
\texttt{segments.shp} & Shapefile & Vectorized polygons with attributes: \texttt{label\_id} (segment ID), \texttt{area\_m2} (area in square meters), \texttt{perimeter\_m} (perimeter in meters). \\
\texttt{segments.gpkg} & GeoPackage & Same content as Shapefile. Generated only when \texttt{save\_geopackage=True}. \\
\texttt{stats.json} & JSON & Processing statistics: total segments, coverage percentage, processing time, per-tile statistics. \\
\bottomrule
\end{tabular}
\end{table}

\subsubsection{Integration with GIS software}

The GeoTIFF and Shapefile/GeoPackage outputs load directly into QGIS, ArcGIS, or any GDAL-compatible software. In Python, the outputs integrate with standard geospatial libraries:

\begin{lstlisting}[style=python]
import geopandas as gpd
import rasterio

# Load vector polygons
segments = gpd.read_file("output/segments.shp")
print(f"Total segments: {len(segments)}")
print(f"Mean area: {segments.area_m2.mean():.1f} m^2")

# Load raster labels
with rasterio.open("output/labels.tif") as src:
    labels = src.read(1)
    print(f"Unique segments: {labels.max()}")
\end{lstlisting}

\subsubsection{Polygon simplification}

The \texttt{simplify\_tolerance} parameter controls the Douglas-Peucker simplification applied to vectorized polygons. The tolerance is specified in map units (typically meters for projected CRS). Table~\ref{tab:simplify} provides guidance for selecting the tolerance.

\begin{table}[htbp]
\centering
\caption{Polygon simplification guidance.}
\label{tab:simplify}
\begin{tabular}{lll}
\toprule
Tolerance & Effect & Use case \\
\midrule
0 & No simplification, all vertices retained & Maximum boundary precision \\
1.0 (default) & Light simplification & General purpose \\
3.0 & Moderate simplification & Faster rendering, smaller files \\
5.0 & Heavy simplification & Overview maps, web display \\
\bottomrule
\end{tabular}
\end{table}


% ----------------------------------------------------------------------------
\subsection{Command-Line Interface}
\label{sec:cli}
% ----------------------------------------------------------------------------

After installation, the \texttt{sam-mosaic} command is available system-wide.

\subsubsection{Basic usage}

\begin{lstlisting}[style=bash]
# Default configuration
sam-mosaic input.tif output/ \
    --checkpoint checkpoints/sam2.1_hiera_large.pt

# With custom tile size for urban scenes
sam-mosaic input.tif output/ \
    --checkpoint checkpoints/sam2.1_hiera_large.pt \
    --tile-size 250 --point-strategy dense_grid

# With configuration file
sam-mosaic input.tif output/ --config config.yaml

# Also generate GeoPackage
sam-mosaic input.tif output/ \
    --checkpoint checkpoints/sam2.1_hiera_large.pt \
    --geopackage
\end{lstlisting}

\subsubsection{All CLI options}

The CLI mirrors the Python API parameters. All parameters listed in Table~\ref{tab:params} are available as command-line flags with hyphenated names (e.g., \texttt{tile\_size} becomes \texttt{-{}-tile-size}). Boolean parameters use \texttt{-{}-no-} prefix to disable (e.g., \texttt{-{}-no-black-mask}).

\subsubsection{YAML configuration file}

For reproducibility, all parameters can be specified in a YAML file:

\begin{lstlisting}[style=yaml]
tile:
  size: 1000
  padding: 50

threshold:
  iou_start: 0.93
  iou_end: 0.60
  stability_start: 0.93
  stability_end: 0.60
  step: 0.01

segmentation:
  points_per_side: 64
  target_coverage: 99.0
  use_black_mask: true
  use_adaptive_threshold: true
  point_strategy: dense_grid
  erosion_iterations: 0

merge:
  merge_strategy: best_match
  min_mask_area: 100
  merge_enclosed_max_area: 500

output:
  simplify_tolerance: 1.0
  save_shapefile: true
  save_geopackage: false
  streaming_mode: auto

sam_checkpoint: checkpoints/sam2.1_hiera_large.pt
\end{lstlisting}

\noindent
Command-line arguments override values from the configuration file, and both override the built-in defaults.


% ----------------------------------------------------------------------------
\subsection{Processing Large Images}
\label{sec:large}
% ----------------------------------------------------------------------------

Remote SAMsing processes images of arbitrary size with constant GPU memory. Three mechanisms enable this:

\subsubsection{Tile-by-tile processing}

Each tile is loaded, segmented, and the result written to the mosaic before the next tile begins. GPU memory holds only one tile at a time. For an $8{,}000 \times 8{,}000$ pixel image at $T = 1{,}000$, this produces 64 tiles processed sequentially.

\subsubsection{Streaming mode}

The \texttt{streaming\_mode} parameter controls how the output label mosaic is stored during processing:
\begin{itemize}
    \item \texttt{"ram"}: The entire mosaic is held in memory as a NumPy array. Fastest option, but requires sufficient RAM for the full image at 32-bit integer resolution (e.g., an $8{,}000 \times 8{,}000$ image requires ${\sim}256$~MB).
    \item \texttt{"disk"}: The mosaic is stored as a memory-mapped file on disk. Enables processing images larger than available RAM at the cost of slower I/O.
    \item \texttt{"auto"} (default): Switches to disk streaming when the mosaic would exceed 30\% of available RAM.
\end{itemize}

\subsubsection{GPU memory management}

SAM2 accumulates GPU memory over repeated inference calls due to internal CUDA state. Remote SAMsing periodically resets the GPU state at configurable intervals (default: every 50 tiles for jobs with 500+ tiles, every 100 tiles otherwise). This reset adds negligible overhead but prevents out-of-memory crashes during long runs.

\subsubsection{Expected processing times}

Table~\ref{tab:times} reports processing times measured on an NVIDIA RTX 4090 (24~GB VRAM) with the Dense Grid point strategy.

\begin{table}[htbp]
\centering
\caption{Processing times for different image sizes and tile configurations. All times measured with Dense Grid on an RTX 4090.}
\label{tab:times}
\begin{tabular}{lrrrl}
\toprule
Image size & Tiles & $T$ & Time & Scene type \\
\midrule
$6{,}000 \times 6{,}000$ & 36 & 1000 & ${\sim}40$~min & Urban (Potsdam, 5~cm) \\
$8{,}000 \times 8{,}000$ & 64 & 1000 & ${\sim}3.6$~h & Urban (BSB, 24~cm) \\
$8{,}000 \times 8{,}000$ & 1024 & 250 & ${\sim}18$~h & Urban (BSB, 24~cm) \\
$10{,}000 \times 10{,}000$ & 100 & 1000 & ${\sim}1.1$~h & Agricultural (Planet, 4.78~m) \\
$36{,}000 \times 54{,}000$ & 1944 & 1000 & ${\sim}12$~h & Urban mosaic (Potsdam) \\
\bottomrule
\end{tabular}
\end{table}

\noindent
Processing time scales approximately linearly with the number of tiles and is dominated by SAM2 inference. Reducing tile size increases the number of tiles quadratically (e.g., $T = 250$ produces $16\times$ more tiles than $T = 1{,}000$).


% ----------------------------------------------------------------------------
\subsection{Practical Recommendations}
\label{sec:recommendations}
% ----------------------------------------------------------------------------

This section provides scene-specific configuration guidance derived from the experimental analysis in the companion paper \citep{carvalho2026jag}.

\subsubsection{Tile size selection}

Tile size determines the spatial scale at which SAM2 operates. The following guidelines apply:

\begin{itemize}
    \item \textbf{Urban scenes at $\leq$~24~cm GSD} (individual cars, narrow walkways visible): Use $T = 250$. This magnifies objects $4\times$ in SAM2's internal resolution, achieving 85\% Det@0.5 on cars and buildings at 24~cm resolution. Processing time is ${\sim}18$~hours for $8{,}000 \times 8{,}000$ pixels.

    \item \textbf{Urban scenes at 5~cm GSD} (objects already large relative to pixel size): Use $T = 1{,}000$. Objects are already well-resolved; reducing tile size fragments large structures into chain-merged mega-segments. Det@0.5 drops from 62\% ($T = 1{,}000$) to 50\% ($T = 250$) on Potsdam at 5~cm.

    \item \textbf{Agricultural scenes at 1--5~m GSD} (large fields, natural vegetation): Use $T = 1{,}000$. Large homogeneous regions benefit from the full context window. BIoU is 0.89 at $T = 1{,}000$ and drops to 0.41 at $T = 250$ on Planet imagery at 4.78~m.

    \item \textbf{Mixed scenes}: Start with $T = 1{,}000$ and inspect the output. If small objects are under-segmented, run again at $T = 500$ or $T = 250$.
\end{itemize}

\subsubsection{Point strategy selection}

Two point strategies are available for subsequent passes (the first pass always uses a uniform grid):

\begin{itemize}
    \item \textbf{Dense Grid} (\texttt{point\_strategy="dense\_grid"}): Regenerates the uniform grid and discards points on already-segmented pixels. Maintains uniform density across all residual regions regardless of their size. Recommended for most use cases.

    \item \textbf{K-means} (\texttt{point\_strategy="kmeans"}): Clusters a fixed number of points in residual areas. Concentrates sampling in large residual regions but may under-sample small fragments. Available for backward compatibility.
\end{itemize}

\subsubsection{Memory considerations}

For images exceeding $20{,}000$ pixels in either dimension:
\begin{itemize}
    \item Set \texttt{streaming\_mode="disk"} to prevent out-of-memory errors during mosaic assembly.
    \item Monitor GPU memory: if CUDA out-of-memory errors occur during SAM2 inference, reduce \texttt{tile\_size} (which also reduces the padded window size).
    \item Set \texttt{save\_geopackage=True} instead of Shapefile for outputs exceeding 2~GB (Shapefile format has a 2~GB size limit).
\end{itemize}


% ============================================================================
\section{Method Validation}
\label{sec:validation}
% ============================================================================

The companion paper \citep{carvalho2026jag} presents a full scientific evaluation of Remote SAMsing on seven scenes from three datasets spanning 5~cm to 4.78~m ground sampling distance, two spectral compositions (natural RGB and MNF false-color), and two landscape types (urban and agricultural). This section summarizes the key validation results; readers are directed to the companion paper for the complete analysis.

\subsection{Coverage}

The multi-pass algorithm achieves 91--98\% coverage at the default tile size ($T = 1{,}000$), compared to 6--27\% for SamGeo2 \citep{wu2023samgeo} (single-pass SAM2 with default settings) and 30--68\% for a single SAM2 pass at the highest quality threshold ($\tau = 0.93$). Reducing tile size to $T = 250$ increases coverage to 97--99\% at the cost of longer processing times. The ablation analysis shows that adaptive threshold decay contributes 3--20 percentage points and black mask focusing contributes 4--11 percentage points to the total coverage.

\subsection{Detection quality}

The greedy oracle evaluation protocol (Det@0.5) measures the fraction of ground truth objects reconstructed with IoU $\geq 0.5$. At the best configuration per scene, Remote SAMsing achieves:
\begin{itemize}
    \item BSB-1 (24~cm, urban): 85.1\% Det@0.5 at $T = 250$, with per-class scores of 95\% for buildings and 93\% for cars.
    \item Potsdam-1 (5~cm, urban): 61.9\% Det@0.5 at $T = 1{,}000$.
    \item Plant23 (4.78~m, agricultural): 95.8\% Det@0.5 at $T = 1{,}000$.
\end{itemize}

\subsection{Boundary quality}

Boundary IoU (BIoU) \citep{cheng2021boundary} measures boundary precision within a narrow band around object edges. Remote SAMsing achieves BIoU 3--8$\times$ higher than SLIC \citep{achanta2012slic} and Felzenszwalb \citep{felzenszwalb2004efficient} baselines across all evaluated scenes.

\subsection{Scalability}

A scalability test on the full Potsdam mosaic ($36{,}000 \times 54{,}000$ pixels, 1.94 billion pixels) produced 124{,}180 segments at 97\% coverage with 81.8\% ASA, confirming that the pipeline operates at production scale. The merge algorithm processed all tile boundaries in under 2 minutes.


% ============================================================================
\section{Limitations}
\label{sec:limitations}
% ============================================================================

\begin{itemize}
    \item \textbf{Processing time.} The multi-pass approach is inherently slower than single-pass segmentation. An $8{,}000 \times 8{,}000$ pixel image at $T = 250$ requires ${\sim}18$ hours on an RTX 4090. Processing time scales quadratically with tile size reduction and linearly with image area.

    \item \textbf{SAM2 dependency.} Remote SAMsing depends on the SAM2.1 Hiera Large checkpoint. The pipeline does not modify SAM2's weights, but it requires a specific PyPI distribution (JinsuaFeito-dev fork) for stable GPU memory during extended runs. Future SAM versions may require adaptation.

    \item \textbf{RGB input only.} SAM2 accepts three-channel RGB input. Remote SAMsing passes any three-band raster to SAM2 (including false-color compositions such as MNF), but single-band (panchromatic), SAR (complex-valued), and thermal imagery have not been tested. Multi-band imagery must be reduced to three channels before processing.

    \item \textbf{No semantic labels.} The pipeline produces unsupervised segments without class labels. Each segment receives a unique integer ID, but the correspondence between segments and land cover classes requires a downstream classification step (e.g., majority-vote labeling in an OBIA framework).

    \item \textbf{Black mask artifacts.} Setting segmented pixels to black creates artificial edges that can affect segment shapes in later passes. While this mechanism is essential for coverage (disabling it reduces coverage by 4--11 pp), segments generated in late passes may have boundaries influenced by the mask rather than by the underlying image content.

    \item \textbf{Tile size sensitivity.} The optimal tile size depends on the scene type and resolution. No automatic selection mechanism exists; the user must select the tile size based on the guidelines in Section~\ref{sec:recommendations} or through trial runs.
\end{itemize}


% ============================================================================
% References
% ============================================================================

\section*{Declaration of Competing Interest}
The authors declare that they have no known competing financial interests or personal relationships that could have appeared to influence the work reported in this paper.

\section*{Data Availability}
The source code is available at \url{https://github.com/osmarluiz/sam-mosaic}. The ISPRS Potsdam dataset is publicly available from the ISPRS 2D Semantic Labeling Contest. The Bras\'{i}lia dataset is publicly available from \citet{carvalho2022panoptic}.

\section*{Acknowledgments}
This study was funded by the Coordination for the Improvement of Higher Education Personnel (CAPES), the National Council for Scientific and Technological Development (CNPq), and the Foundation for Research Support of the Federal District (FAPDF).

\section*{CRediT authorship contribution statement}
\textbf{Osmar Luiz Ferreira de Carvalho}: Conceptualization, Methodology, Software, Investigation, Writing -- original draft.
\textbf{Osmar Ab\'{i}lio de Carvalho J\'{u}nior}: Supervision, Writing -- review \& editing, Funding acquisition.
\textbf{Anesmar Olino de Albuquerque}: Investigation, Data curation, Writing -- review \& editing.
\textbf{Daniel Guerreiro e Silva}: Writing -- review \& editing, Supervision.

\bibliographystyle{elsarticle-harv}

\begin{thebibliography}{25}

\bibitem[{Achanta et~al.(2012)}]{achanta2012slic}
Achanta, R., Shaji, A., Smith, K., Lucchi, A., Fua, P., S\"{u}sstrunk, S., 2012. SLIC Superpixels Compared to State-of-the-Art Superpixel Methods. IEEE Transactions on Pattern Analysis and Machine Intelligence 34~(11), 2274--2282.

\bibitem[{Carvalho et~al.(2026)}]{carvalho2026jag}
de~Carvalho, O.L.F., de~Carvalho~J\'{u}nior, O.A., de~Albuquerque, A.O., Guerreiro~e~Silva, D., 2026. Remote SAMsing: From Segment Anything to Segment Everything. International Journal of Applied Earth Observation and Geoinformation (companion paper).

\bibitem[{Carvalho et~al.(2022)}]{carvalho2022panoptic}
de~Carvalho, O.L.F., de~Carvalho~J\'{u}nior, O.A., Silva, C.R., de~Albuquerque, A.O., Santana, N.C., Borges, D.L., Gomes, R.A.T., Guimar\~{a}es, R.F., 2022. Panoptic Segmentation Meets Remote Sensing. Remote Sensing 14~(4), 965.

\bibitem[{Cheng et~al.(2021)}]{cheng2021boundary}
Cheng, B., Girshick, R., Doll\'{a}r, P., Berg, A.C., Kirillov, A., 2021. Boundary IoU: Improving Object-Centric Image Segmentation Evaluation. In: Proceedings of the IEEE/CVF Conference on Computer Vision and Pattern Recognition, pp.~15334--15342.

\bibitem[{Felzenszwalb and Huttenlocher(2004)}]{felzenszwalb2004efficient}
Felzenszwalb, P.F., Huttenlocher, D.P., 2004. Efficient Graph-Based Image Segmentation. International Journal of Computer Vision 59~(2), 167--181.

\bibitem[{Kirillov et~al.(2023)}]{kirillov2023segment}
Kirillov, A., Mintun, E., Ravi, N., Mao, H., Rolland, C., Gustafson, L., Xiao, T., Whitehead, S., Berg, A.C., Lo, W.-Y., et~al., 2023. Segment Anything. In: Proceedings of the IEEE/CVF International Conference on Computer Vision, pp.~4015--4026.

\bibitem[{Osco et~al.(2023)}]{osco2023sam_rs_review}
Osco, L.P., Wu, Q., de~Lemos, E.L., Gon\c{c}alves, W.N., Ramos, A.P.M., Li, J., Marcato~Junior, J., 2023. The Segment Anything Model (SAM) for Remote Sensing Applications: From Zero to One Shot. International Journal of Applied Earth Observation and Geoinformation 124, 103540.

\bibitem[{Ravi et~al.(2024)}]{ravi2024sam2}
Ravi, N., Gabeur, V., Hu, Y.-T., Hu, R., Ryali, C., Ma, T., Khedr, H., R\"{a}dle, R., Rolland, C., Gustafson, L., et~al., 2024. SAM 2: Segment Anything in Images and Videos. arXiv preprint arXiv:2408.00714.

\bibitem[{Stutz et~al.(2018)}]{stutz2018superpixels}
Stutz, D., Hermans, A., Leibe, B., 2018. Superpixels: An Evaluation of the State of the Art. Computer Vision and Image Understanding 166, 1--27.

\bibitem[{Tarjan(1975)}]{tarjan1975union}
Tarjan, R.E., 1975. Efficiency of a Good But Not Linear Set Union Algorithm. Journal of the ACM 22~(2), 215--225.

\bibitem[{Wu and Osco(2023)}]{wu2023samgeo}
Wu, Q., Osco, L.P., 2023. samgeo: A Python package for segmenting geospatial data with the Segment Anything Model (SAM). Journal of Open Source Software 8~(89), 5663.

\end{thebibliography}

\end{document}
